%% Resumé
\makeatletter
\newenvironment{abstracts}{%
  \ifx\maketitle\relax
    \ClassWarning{\@classname}{Abstract should precede
      \protect\maketitle\space in AMS document classes; reported}%
  \fi
  \global\setbox\abstractbox=\vtop \bgroup
    \normalfont\Small
    \list{}{\labelwidth\z@
      \leftmargin3pc \rightmargin\leftmargin
      \listparindent\normalparindent \itemindent\z@
      \parsep\z@ \@plus\p@
      \let\fullwidthdisplay\relax
      \itemsep\medskipamount
    }%
}{%
  \endlist\egroup
  \ifx\@setabstract\relax \@setabstracta \fi
}

\newcommand{\abstractin}[1]{%
  \otherlanguage{#1}%
  \item[\hskip\labelsep\scshape\abstractname.]%
}
\makeatother


%% Gestion des théorèmes & Cie.

\swapnumbers % Le numéro des théorèmes s'affiche avant le mot "Théorème".
\theoremstyle{definition} % Pour tout ce qui ressemble à des définitions : définitions, notations...
\newtheorem{Def}{Definition}[section] % Définitions.
\theoremstyle{plain} % Pour tous ce qui ressemble à des théorèmes.
\newtheorem{Pro}[Def]{Proposition} % Propositions.
\newtheorem{Lem}[Def]{Lemma} % Lemmes.
\newtheorem{The}[Def]{Theorem} % Théorèmes.
\newtheorem{Cor}[Def]{Corollary} % Corollaires.
\newtheorem{Con}[Def]{Conjecture} % Corollaires.
%
%\newtheorem{Defa}[Def]{Definition} % Propositions.
%\newtheorem{Proa}[Def]{Proposition} % Propositions.
%\newtheorem{Lema}[Def]{Lemma} % Lemmes.
%\newtheorem{Thea}[Def]{Theorem} % Théorèmes.
%\newtheorem{Cora}[Def]{Corollary} % Corollaires.
%\newtheorem{Cona}[Def]{Conjecture} % Corollaires.

%
%
%\newtheorem*{The*}{Théorème}
%\newtheorem*{Def*}{Définition}
%\newtheorem*{Pro*}{Proposition}
\newtheorem*{The*}{Theorem}
\newtheorem*{Def*}{Definition} % Propositions.
\newtheorem*{Pro*}{Proposition}

\theoremstyle{remark} % Pour tout ce qui ressemble à des remarques.
\newtheorem{Exe}[Def]{Example} % Exemples.
\newtheorem{Rem}[Def]{Remark} % Remarques.

%
%
%\newtheorem*{Rem*}{Remarque}
\newtheorem*{Rem*}{Remark}

%
%
%%%Changement numérotation annexe
%\newtheorem{DefA}{Définition}[section] % Définitions.
%\theoremstyle{plain} % Pour tous ce qui ressemble à des théorèmes.
%\newtheorem{ProA}[DefA]{Proposition} % Propositions.
%\newtheorem{LemA}[DefA]{Lemme} % Lemmes.
%\newtheorem{TheA}[DefA]{Théorème} % Théorèmes.
%\newtheorem{CorA}[DefA]{Corollaire} % Corollaires.
%\newtheorem{ConA}[DefA]{Conjecture} % Corollaires.
%


%%Chapitre sans saut de page
\makeatletter
\newcommand\chapterpage{\clearpage
                    \thispagestyle{plain}%
                    \global\@topnum\z@
                    \@afterindentfalse
                    \secdef\@chapter\@schapter}
\makeatother


%% Groupes algébriques et duaux
\makeatletter
\def\parenthese#1{\@ifnextchar_{(#1)}{#1}}
\makeatother

\newcommand{\gp}[1]{#1}
\NewDocumentCommand{\algPoints}{m o}
{
	\IfNoValueTF {#2}
		{#1}
    	{\parenthese{#1(#2)}} 	
}
\NewDocumentCommand{\gpalg}{ s m o}{
	\IfBooleanTF#1
    	{\algPoints{\widehat{\textbf{#2}}}[#3]}
    	{\algPoints{\textbf{#2}}[#3]}
}
\NewDocumentCommand{\gpfini}{ s m}{
	\IfBooleanTF#1
    	{\textsf{#2}^{*}}
    	{\textsf{#2}}
}
\NewDocumentCommand{\gpfinialg}{s m o}{
	\IfBooleanTF#1
    	{\algPoints{\textbf{\textsf{#2}}^{*}}[#3]}
    	{\algPoints{\textbf{\textsf{#2}}}[#3]}
}
\NewDocumentCommand{\Ldual}{m o}{\algPoints{ {}^{L}\gpalg{#1}}[#2]}
\NewDocumentCommand{\quotred}{s m m}{
	\parenthese{
		\IfBooleanTF#1
			{\overline{\gpfini{#2}}_{#3}^{*} }
			{\overline{\gpfini{#2}}_{#3} }
	}
}
\NewDocumentCommand{\quotredalg}{s m m o}{
	\parenthese{ 
		\IfBooleanTF#1
			{\algPoints{\overline{\gpfinialg{#2}}_{#3}^{*}}[#4]}
			{\algPoints{\overline{\gpfinialg{#2}}_{#3}}[#4]}
	}
}
\newcommand{\G}{\gp{G}}


\NewDocumentCommand{\cent}{s m m}{ 
	\IfBooleanTF#1
		{C_{#3}(#2)^{\circ}}
		{C_{#3}(#2)}
}



\DeclareDocumentCommand{\ss}{s m}{
	\IfBooleanTF#1
		{#2_{ss,\Lambda}}
		{#2_{ss}}
}


%% Groupes classiques
\newcommand{\ZnZ}[1]{\mathbb{Z}/#1\mathbb{Z}} %Z/nZ
\DeclareMathOperator{\GL}{GL}
\DeclareMathOperator{\SL}{SL}
\DeclareMathOperator{\PGL}{PGL}
\DeclareMathOperator{\SO}{SO}
\DeclareMathOperator{\Sp}{Sp}
\DeclareMathOperator{\Ortho}{O}
\DeclareMathOperator{\Un}{U}
\NewDocumentCommand{\gl}{m o}{\algPoints{\GL_{#1}}[#2]}  
\DeclareDocumentCommand{\sl}{m o}{\algPoints{\SL_{#1}}[#2]}  
\NewDocumentCommand{\pgl}{m o}{\algPoints{\PGL_{#1}}[#2]}  
\DeclareDocumentCommand{\so}{s m o}{
	\IfBooleanTF#1
		{\algPoints{\SO_{#2}^{*}}[#3]}
		{\algPoints{\SO_{#2}}[#3]}
}
\DeclareDocumentCommand{\sp}{m o}{\algPoints{\Sp_{#1}}[#2]} 
\NewDocumentCommand{\on}{m o}{\algPoints{\Ortho_{#1}}[#2]} 
\NewDocumentCommand{\un}{m o o}{\algPoints{\Un_{#1}}[#2/#3]} 




%%% Groupes de Weyl

\newcommand{\WfM}[1]{W_{#1}}
\newcommand{\WaM}[1]{W_{#1}^a}
\newcommand{\WeM}[1]{\tilde{W}_{#1}}

\NewDocumentCommand{\Wf}{o}{
	\IfNoValueTF {#1}
	{W}
	{W(\gpalg{#1})}
}
\NewDocumentCommand{\Wa}{o}{
	\IfNoValueTF {#1}
		{W^a}
		{W^a(\gpalg{#1})}
}
\NewDocumentCommand{\We}{o}{
	\IfNoValueTF {#1}
		{\tilde{W}}
		{\tilde{W}(\gpalg{#1})}
}
\NewDocumentCommand{\Wx}{m o}{
	\IfNoValueTF {#2}
		{W_{#1}}
		{W_{#1}(\gpalg{#2})}
}

\NewDocumentCommand{\Wxd}{m o}{
	\IfNoValueTF {#2}
		{W_{#1}^{\dagger}}
		{W_{#1}^{\dagger}(#2)}
}
\NewDocumentCommand{\WStf}{s}{
	\IfBooleanTF#1
		{W(\gpalg{S},\theta)}
		{W^{\circ}(\gpalg{S},\theta)}
}
\NewDocumentCommand{\WStxcomplet}{m m m}{W^{\circ}_{#1}(#2,#3)}
\NewDocumentCommand{\WStx}{m}{\WStxcomplet{#1}{\gpalg{S}}{\theta}}
\NewDocumentCommand{\WStecomplet}{s m m}{
	\IfBooleanTF#1
		{\tilde{W}(#2,#3)}
		{\tilde{W}^{\circ}(#2,#3)}
}
\NewDocumentCommand{\WSte}{s}{
	\IfBooleanTF#1
		{\WStecomplet*{\gpalg{S}}{\theta}}
		{\WStecomplet{\gpalg{S}}{\theta}}
}
\newcommand{\WStacomplet}[2]{W^{a}(#1,#2)}
\newcommand{\WSta}{\WStacomplet{\gpalg{S}}{\theta}}

\NewDocumentCommand{\WsM}{s m m}{
	\IfBooleanTF#1
		{W_{#3,#2}}
		{W^{\circ}_{#3,#2}}
}

\NewDocumentCommand{\WSteM}{s m}{
	\IfBooleanTF#1
		{\tilde{W}_{#2}(\gpalg{S},\theta)}
		{\tilde{W}^{\circ}_{#2}(\gpalg{S},\theta)}
}
\newcommand{\WStaM}[1]{W^{a}_{#1}(\gpalg{S},\theta)}
\NewDocumentCommand{\WStfM}{s m}{
	\IfBooleanTF#1
		{W_{#2}(\gpalg{S},\theta)}
		{W^{\circ}_{#2}(\gpalg{S},\theta)}
}
\NewDocumentCommand{\Ws}{s m o}{
	\IfBooleanTF#1
	{
		\IfNoValueTF {#3}
			{W_{#2}}
			{W_{#2}(#3)}
	}
	{
		\IfNoValueTF {#3}
			{W^{\circ}_{#2}}
			{W^{\circ}_{#2}(#3)}
	}
}

\NewDocumentCommand{\NStecomplet}{s m m}{
	\IfBooleanTF#1
		{N(#2,#3)}
		{N^{\circ}(#2,#3)}
}
\NewDocumentCommand{\NSte}{s}{
	\IfBooleanTF#1
		{\NStecomplet*{\gpalg{S}}{\theta}}
		{\NStecomplet{\gpalg{S}}{\theta}}
}
\newcommand{\NStacomplet}[2]{N^a(#1,#2)}
\newcommand{\NSta}{\NStacomplet{\gpalg{S}}{\theta}}

\newcommand{\Wap}{W'^a}
%% Bruhat-Tits

\DeclareMathOperator{\BT}{BT}
\NewDocumentCommand{\bt}{o}{
	\IfNoValueTF{#1}
		{\BT}
		{\BT(#1)}
}
\NewDocumentCommand{\bts}{o}{
	\IfNoValueTF{#1}
		{\BT_{0}}
		{\BT_{0}(#1)}
}
\NewDocumentCommand{\bte}{o}{
	\IfNoValueTF{#1}
		{\BT^{e}}
		{\BT^{e}(#1)}
}
\NewDocumentCommand{\btes}{o}{
	\IfNoValueTF{#1}
		{\BT^{e}_{0}}
		{\BT^{e}_{0}(#1)}
}
\newcommand{\app}[2]{\mathcal{A}(\gpalg{#1},#2)}
\newcommand{\appe}[2]{\mathcal{A}^e(\gpalg{#1},#2)}

\newcommand{\fix}[2]{#1_{#2}}
\newcommand{\para}[2]{#1_{#2}^{\circ}}
\newcommand{\radpara}[2]{#1_{#2}^{+}}
\newcommand{\paranr}[2]{#1_{#2}^{\circ,nr}}
\newcommand{\radparanr}[2]{#1_{#2}^{+,nr}}


%% Systèmes de conjugaisons

\NewDocumentCommand{\systConj}{s m}{
	\IfBooleanTF#1
		{\fc{#2}=(\fc{#2}_{\sigma})_{\sigma \in \bt}}
		{\fc{#2}}
}
\NewDocumentCommand{\systConjx}{m m}{\fc{#1}_{#2}}
\NewDocumentCommand{\systEng}{m m}{\fc{S}_{(#1,#2)}}
\NewDocumentCommand{\systEngx}{m m m}{\fc{S}_{(#1,#2),#3}}
\newcommand{\compl}[1]{{}^{c}#1}
\NewDocumentCommand{\systConjC}{m}{\fc{#1}=(\fc{#1}_{\sigma})_{\sigma \leq C}}
\NewDocumentCommand{\systConjCpro}{m}{\fc{#1}=(\fc{#1}_{\sigma})_{\sigma \leq C}}

\newcommand{\fc}[1]{\dofc#1}
\newcommand{\dofc}[1]{\mathcal{#1}}

\newcommand{\clotco}[1]{\overline{\fc{#1}}}
\newcommand{\SStheta}{\systConj{S}_{(\gpalg{S},\theta)}}
\newcommand{\SA}[1]{\systConj{S}_{#1}}

\NewDocumentCommand{\transss}{s m m}{
	\IfBooleanTF#1
		{\varphi^{*}_{#2,#3}}
		{\varphi_{#2,#3}}
}

\newcommand{\simJss}{\sim}
\newcommand{\simJsst}{\sim'}


%%Ensembles

\newcommand{\ensSystConj}{\mathscr{S}_{\ld}}
\newcommand{\ensSystConjc}{\ensSystConj^{c}}
\newcommand{\Smin}{\ensSystConj^{\text{min}}}
\newcommand{\ensSystConjcM}{\mathscr{S}_{M,\ld}^{c}}
\newcommand{\pairesStmin}{\mathscr{P}_{\ld}^{\text{min}}}
\newcommand{\triplets}{\mathscr{T}_{\ld}}
\newcommand{\tripletsm}{\mathscr{T}^{m}_{\ld}}
\newcommand{\tripletse}{\mathscr{T}^{e}_{\ld}}
\newcommand{\Jss}{\mathscr{C}_{\ld}}
\newcommand{\parties}[1]{\mathcal{P}(#1)}

\NewDocumentCommand{\pairesSt}{s}{
	\IfBooleanTF#1
	{\mathscr{P}^{*}_{\ld}}
	{\mathscr{P}_{\ld}}
}
\newcommand{\classer}[2]{[\gpalg{#1},#2]_r}
\newcommand{\clr}{\classer{S}{\theta}}

%% Paramètres de Langlands
\NewDocumentCommand{\Lparametre}{m o}{\Phi(\Ldual{#1}[#2])}
\NewDocumentCommand{\Lparam}{m m o}{\Phi(#1,\Ldual{#2}[#3])}
\NewDocumentCommand{\Lparamm}{m m o}{\Phi_m(#1,\Ldual{#2}[#3])}
\NewDocumentCommand{\Lp}{m o}{\Lparam{#1}{\G}[#2]}
\NewDocumentCommand{\Lpm}{m o}{\Lparamm{#1}{\G}[#2]}
\NewDocumentCommand{\LparamB}{m m o}{\widetilde{\Phi}(#1,\Ldual{#2}[#3])}
\NewDocumentCommand{\LparamBm}{m m o}{\widetilde{\Phi}_m(#1,\Ldual{#2}[#3])}
\NewDocumentCommand{\Lpb}{m o}{\LparamB{#1}{\G}[#2]}
\NewDocumentCommand{\Lpbm}{m o}{\LparamBm{#1}{\G}[#2]}


%%Nombres l-adics
\newcommand{\ladic}[1]{\overline{\mathbb{#1}}_{\lprime}}
\newcommand{\Ql}{\ladic{Q}}
\newcommand{\Zl}{\ladic{Z}}
\newcommand{\Fl}{\ladic{F}}


%%Représentations
\DeclareMathOperator{\Rep}{Rep}
\DeclareMathOperator{\Irr}{Irr}
\NewDocumentCommand{\rep}{ O{} O{} m }{\Rep_{#1}^{#2}(#3)}
\NewDocumentCommand{\indPara}{m m o}{
	\IfNoValueTF{#3}
		{i_{#1}^{#2}}
		{i_{#1}^{#2}(#3)}
}
\NewDocumentCommand{\resPara}{m m o}{
	\IfNoValueTF{#3}
		{r_{#1}^{#2}}
		{r_{#1}^{#2}(#3)}
}
\NewDocumentCommand{\ip}{ o}{\indPara{\gp{P}}{\G}[#1]}
\NewDocumentCommand{\rp}{ o}{\resPara{\gp{P}}{\G}[#1]}


%%Operateurs
\DeclareMathOperator{\Aut}{Aut}
\DeclareMathOperator{\Out}{Out}
\DeclareMathOperator{\Int}{Int}
\DeclareMathOperator{\Gal}{Gal}
\DeclareMathOperator{\Ad}{Ad}
\DeclareMathOperator{\Tr}{Tr}
\DeclareMathOperator{\Hom}{Hom}
\DeclareMathOperator{\Res}{Res}



%% Mise en forme
\newcommand{\sautintro}{\bigskip}

\newcommand{\tosim}{\overset{\sim}{\to}}

%%Notations
\newcommand{\lprime}{\ell}
\newcommand{\ld}{\Lambda}
\newcommand{\ldb}{\overline{\ld}}
\newcommand{\hecke}[2]{\mathcal{H}_{#2}(#1)}

\newcommand{\iner}{I_{\kk}}
\newcommand{\inerl}{\iner^{(\lprime)}}
\newcommand{\Ild}{\iner^{\ld}}
\newcommand{\weil}{W_{\kk}}
\newcommand{\inersauv}{P_{\kk}}

\newcommand{\dl}[2]{\mathcal{E}(#1,#2)}

\newcommand{\sd}{s_{t}}
\newcommand{\cd}{[\varphi_{t}]}

\newcommand{\kk}{F}
\newcommand{\knr}{F^{nr}}
\newcommand{\kalg}{\overline{F}}
\newcommand{\res}{k}
\newcommand{\resalg}{\overline{k}}

\newcommand{\fr}{\mathsf{F}}




%%%%%%%%%%%%%%%%%%%%%%%%%%% NOuveaux %%%%%%%%%%%%%%%%%%%%%%%%%%%%%%%%
\newcommand{\Gphisigma}[2]{\gpalg{G}_{#1,#2}}
\newcommand{\Gps}{\Gphisigma{\phi}{\sigma}}
\newcommand{\Gpsc}{\Gps^{\circ}}
\newcommand{\Gpsnr}{G_{\phi,\sigma}^{nr}}
\newcommand{\Gpscnr}{G_{\phi,\sigma}^{\circ,nr}}
\newcommand{\Mpsc}{\gpalg{M}_{\phi_M,\sigma_M}^{\circ}}
\newcommand{\Mpscnr}{M_{\phi_M,\sigma_M}^{\circ,nr}}
\newcommand{\pairesStw}[1]{\mathscr{P}^{#1}_{\ld}}
\newcommand{\der}[1]{\mathcal{D}(#1)}
\newcommand{\hps}{h_{\phi,\sigma}}




%%%%%%%%%%%%%%%%%Nouveaux these$$$$$$$$$$$$$$$$$$$$
\newcommand{\modG}[1]{\delta_{#1}}
\DeclareMathOperator{\Ind}{Ind}
\newcommand{\bern}[1]{\mathcal{B}(#1)}
\DeclareMathOperator{\Coef}{Coef}
\newcommand{\coef}{\Coef_{e}(G,R)}
\newcommand{\coefG}[1]{\Coef_{e}(#1,R)}


\newcommand{\TG}{\mathcal{T}(G)}

\newcommand{\ordpol}[1]{P_{\gpfini{#1}}}
\newcommand{\card}[1]{|#1|}
\newcommand{\cycl}[1]{\Phi_{#1}}

%\DeclareMathOperator{\Rr}{{}^{*}R}
%\newcommand{\resdl}[3]{\Rr_{\gpfinialg{#1} \subseteq \gpfinialg{#2}}^{\gpfinialg{#3}}}
%\newcommand{\rdl}{\resdl{L}{P}{G}}

\newcommand{\resdl}[3]{{}^{*}\mathcal{R}_{\gpfinialg{#1} \subset \gpfinialg{#2}}^{\gpfinialg{#3}}}
\newcommand{\rdl}{\resdl{L}{P}{G}}


%\DeclareMathOperator{\Ri}{R}
%\newcommand{\inddl}[3]{\Ri_{\gpfinialg{#1} \subseteq \gpfinialg{#2}}^{\gpfinialg{#3}}}
%\newcommand{\idl}{\inddl{L}{P}{G}}

\newcommand{\inddl}[3]{\mathcal{R}_{\gpfinialg{#1} \subset \gpfinialg{#2}}^{\gpfinialg{#3}}}
\newcommand{\idl}{\inddl{L}{P}{G}}

\DeclareMathOperator{\Abirr}{Ab_{\lprime}Irr}

\newcommand{\dldcusp}[2]{\dl{\gpfini{G}}{(\gpfinialg{#1},#2)}}

\DeclareMathOperator{\defect}{defect}
\DeclareMathOperator{\rank}{rank}
\newcommand{\quotpara}[2]{\gpfini{P}_{#1,#2}}
\newcommand{\quotunip}[2]{\gpfini{U}_{#1,#2}}

\newcommand{\kgdc}[2]{k(#1,#2)}
\newcommand{\kgd}{\kgdc{\gpfini{G}}{d}}


\newcommand{\Fq}[1]{\mathbb{F}_{q^{#1}}}
\newcommand{\Tc}{\gpfini{T}_c}
\newcommand{\TGu}{\mathcal{T}^{1}(G)}

\newcommand{\dll}[2]{\mathcal{E}_{\lprime}(#1,#2)}

\newcommand{\rl}{r_{\lprime}}
\newcommand{\zG}[1]{Z_{#1}(G)}
\DeclareMathOperator{\ind}{ind}

\DeclareMathOperator{\Gr}{Gr}
\DeclareMathOperator{\im}{Im}

\newcommand{\TGl}{\mathcal{T}_{\lprime}^{1}(G)}

\DeclareMathOperator{\End}{End}
\DeclareMathOperator{\Ker}{Ker}
\DeclareMathOperator{\Id}{Id}